% Copyright Ben Paul Wise. All Rights Reserved.
% ----------------------------------------------
% Part IV -- Mathematical foundations (a regular chapter, for OR
% professionals). Drafted second per the approved order. Proof standard:
% cite, and verify the cited result's conditions against our setting; do not
% reprove. (A Detailed Proofs appendix may be added later from the saved
% reconstruction notes.) Sources: the engine headers,
% network/doc/{formulation,reduction}.md (transcribed, not re-derived), the
% Maxima scripts, and the run evidence of Chapter 3.
% ----------------------------------------------
\chapter{Mathematical Foundations}
\label{ch:math}

This chapter gives the mathematics behind the engines and the three
application problems, at the level an operations-research professional needs
to audit the library's claims: each algorithm is stated as implemented, its
convergence theory is \emph{cited} with the cited result's hypotheses
checked against our setting, and every departure from a source is called
out. Full re-derivations are deliberately omitted. Where our own algebra
goes beyond the literature (the structured flow factory, the game
formulations, the alternating chain), the formulas on which the results
depend were verified symbolically with Maxima \cite{maxima} --- a computer
algebra system in continuous development and improvement since its origins
as MIT's Macsyma \cite{macsyma} --- and the verification scripts are named
where the results are used, so a reader can re-run every check rather than
take any formula on trust.

\section{Problem statement and geometry}
\label{sec:a-problem}

Throughout, the target problem is the \emph{mixed complementarity problem}
(MCP): find $z = (x, y) \in \mathbb{R}^{n} \times \mathbb{R}^{m}$ with
\begin{equation}
  H(x, y) = 0, \qquad 0 \le G(x, y) \perpc y \ge 0,
  \label{eq:mcp}
\end{equation}
equivalently the variational inequality VI$(F, K)$ --- find $z \in K$ with
$F(z)^\top (w - z) \ge 0$ for all $w \in K$ --- over the polyhedral cone
$\K$ with $F = (H, G)$. The affine case $F(z) = Mz + q$ is the mixed linear
complementarity problem (LCP); $n = 0$ gives the pure (N)CP.

\paragraph{The natural map.} With $P_K$ the Euclidean projection onto $K$,
the natural residual
\begin{equation}
  r(z) \;=\; z - P_K\bigl(z - F(z)\bigr)
  \label{eq:natres}
\end{equation}
vanishes exactly at solutions of \eqref{eq:mcp} (Facchinei--Pang's standard
equivalence; for our $K$, $r$ splits into $r_x = H$ on the free block and
$r_i = \min(y_i, G_i)$ componentwise on the orthant block). The library's
universal convergence measure is the \emph{squared} norm $\|r(z)\|^2$;
every engine terminates and reports on it, whatever quantity it steers by
internally, so cross-engine comparisons are well-defined
(Section~\ref{sec:a-measure}).

\paragraph{Problem classes.} $F$ is \emph{monotone} when
$(F(z) - F(w))^\top (z - w) \ge 0$ (for affine $F$: $M$ positive
semidefinite --- the KKT map of a convex QP is the canonical example);
\emph{pseudomonotone} when $F(w)^\top(z - w) \ge 0$ implies
$F(z)^\top(z - w) \ge 0$ (strictly weaker); and otherwise general
(``nonmonotone''). The three application problems occupy three points of
this hierarchy: the flow-planning LCP is monotone by construction
(Section~\ref{sec:a-flowqp}); the influence game is nonmonotone but mild
(its nonlinearity at the shipped setting is a single rational term); the
deployment game is nonmonotone \emph{by design}, with mixed curvature
(Section~\ref{sec:a-ppd}).

\paragraph{Degeneracy.} A solution $z^\ast$ is \emph{nondegenerate} when
strict complementarity holds ($y_i^\ast + G_i^\ast > 0$ for every $i$) and
suitable Jacobian blocks are nonsingular; failures of either produce the
phenomena that drive this library's design: non-isolated solution sets
(wide optimal faces of the QP under near-tied costs), singular semismooth
Newton matrices (Munson et al.\ report singular factorizations in roughly
$12\%$ of iterations across MCPLIB \cite{munson2001}), and
projection-contraction rates approaching one
(Section~\ref{sec:a-pc}).

\section{Projection-contraction methods}
\label{sec:a-pc}

\paragraph{He's method (\code{bsHe94b}).} He's 1994 projection-contraction
method \cite{he1994}, in the paper's equation-(16) search-direction form,
solves the affine VI over any closed convex $K$ reachable by projection. It
fixes the metric $G = M + I$ once, factors it once, and each iteration
computes the projection residual, solves one system against the stored
factorization, and takes a relaxed step ($\gamma \in (0,2)$, default $1.6$).
The contraction argument --- $\|z_{k+1} - z^\ast\|_G^2 \le
\|z_k - z^\ast\|_G^2 - c\,\|r(z_k)\|^2$ --- requires $M$ positive
semidefinite; we verify that hypothesis per problem (for the flow LCP it
holds by construction). Per-iteration cost: one back-substitution and two
projections.

\paragraph{Han's method (\code{dHan06}).} Han's 2006 self-adaptive variant
\cite{han2006} replaces the fixed metric with $(I + \beta_k M)$, adapting
$\beta_k$ by the paper's rule (its eq.~11) under a $\tau$-schedule.
Convergence (Han's Thm~2.4) again assumes monotonicity. Two implementation
notes. First, the shipped default $\beta_0 = 1$ is a deliberate departure
from the reference implementation's $0.5$: it makes the first iteration's
metric equal He's proven-contractive $M + I$, so the method starts stably
even when a Josephy--Newton linearization is indefinite; Han's theory
permits any $\beta_0 > 0$. Second, because $\beta_k$ moves, the method
refactors \emph{every} iteration --- observed at roughly an order of
magnitude more wall time than He's method at equal iteration counts.

\paragraph{Why both stall on degenerate faces.} The contraction constant
degrades with the conditioning of the problem restricted to the optimal
face. Near-tied costs (the banded geometry of
Section~\ref{sec:a-flowqp}) produce wide, nearly flat optimal faces on
which the quadratic growth of the underlying QP objective is almost zero in
many directions; the per-step reduction factor then approaches one, and the
iteration count grows by orders of magnitude, with no warning other than
the slowness itself. This is a property of the method class, not of a
particular tuning: it is the observed 150{,}000-iteration cap on the
200-node banded instance (Chapter~\ref{ch:testing}), and the reason the
interior-point engine exists. We state this as a mechanism rather than a
sharp theorem; the quantitative form would need the face's geometry
explicitly.

\section{Solodov--Svaiter hyperplane projection}
\label{sec:a-ss}

The double-projection method of Solodov and Svaiter \cite{solodov1999}
computes $r = x - P_K(x - \mu F(x))$, backtracks along $r$ until the
Armijo-type condition $\langle F(y), r\rangle \ge (\sigma/\mu)\|r\|^2$
holds at a trial point $y$, and projects the iterate onto the separating
hyperplane $\{z : \langle F(y), z - y\rangle = 0\}$, then back onto $K$.
Its convergence theorem --- global convergence for \emph{continuous
pseudomonotone} $F$, with no Lipschitz constant and nothing factored ---
has the weakest hypotheses of any engine in the library, and its
matrix-free iteration costs only products with $M$.

The library's contribution is a sharpened understanding of its \emph{rate}
in a regime the paper does not emphasize: whenever $\|F(x^\ast)\| > 0$ (the
generic case for a VI whose solution is on the boundary of $K$, e.g.\
linear objectives over an ellipsoid), the hyperplane step length scales
like $\|r\|^2 / \|F\|^2$, and a Fej\'er argument then yields only an
$O(1/\sqrt{k})$ tail on the squared residual. Empirically the calibration
is exact (the ellipsoid experiment of Chapter~\ref{ch:testing}: two Newton-
metric engines converge in 35 and 101 iterations where the hyperplane
method is still at $10^{-3}$ after half a million). The design consequence
is reflected in the API: the method is offered as a \emph{globalizer} ---
phase one of \code{chainedSolodovHe}, handing a warm start to He's
contraction --- never as a finisher, and its own acceptance tests assert
only globalization-grade accuracy.

\section{Forward-backward splitting}
\label{sec:a-fbs}

The two-projection forward-backward splitting scheme of He, Yuan, and
Zhang \cite{hyz2004} --- $\rho = P_K(x - \gamma F(x))$, then
$x^{+} = P_K(\rho + \gamma(F(x) - F(\rho)))$ --- solves the general
monotone Lipschitz VI with two field evaluations and two projections per
iteration and nothing factored; the implementation adapts $\gamma$ from a
smoothed secant estimate of the Lipschitz constant, preserving the rule of
the author's earlier \code{pmedemo} implementation. Its place in the
catalog lies between the affine projection engines and the Newton-class
solvers: it addresses the \emph{nonlinear} problem directly, like the
semismooth solver, but Jacobian-free and over any projector-defined $K$
like the hyperplane method, at a linear rate. Historically it is the
engine of the PME reference implementation, where it solved the
strategic-allocation game of Section~\ref{sec:a-saoe} --- thousands of
iterations where the Newton-class engines take tens, but requiring nothing
beyond evaluations. The library port adds two documented safeguards beyond
the 2004 recipe: a conservative sampled-secant preflight for the initial
step, and Tseng-style step backtracking (halve $\gamma$ until
$\gamma\|F(\rho) - F(x)\| \le 0.9\|\rho - x\|$), which extends robustness
to fields that are only \emph{locally} Lipschitz --- e.g.\ cubic maps over
an unbounded $K$, where no start-point estimate can be globally valid. The
bare recipe sufficed historically because SAOE-class fields are unusually
well behaved: every slope is divided by the total committed effort, so the
field saturates and its Lipschitz constant is nearly uniform. A script
validated on one problem and a library engine required to behave correctly
on any problem satisfy different contracts. On such well-behaved fields
neither safeguard has any effect (the backtracking test costs nothing
unless it fails).

\section{The Josephy--Newton driver}
\label{sec:a-jn}

For the smooth nonlinear VI, \code{solveVI} implements the Josephy--Newton
scheme \cite{josephy1979}: at $z_k$, solve the affine VI with
$M = J(z_k)$, $q = F(z_k) - J(z_k) z_k$ over the \emph{same} $K$, then damp
the step $z_{k+1} - z_k$ by an Armijo backtracking search on the natural
merit $\tfrac12\|r(z)\|^2$. Three implementation decisions matter:

\begin{itemize}
\item \textbf{$K$ is never eliminated.} The subproblem keeps the full mixed
  structure through the projector; there is no Schur complement on the free
  block. This keeps the driver correct for \emph{any} projector-defined
  $K$ (the ellipsoid tests exercise exactly this) at the cost of a larger
  affine subproblem.
\item \textbf{Damping is necessary.} At a non-smooth complementarity
  solution the undamped iteration exhibits a period-2 oscillation
  (overshoot across the kink, get sent back); the value-form Armijo
  condition $\varphi(\alpha) \le (1 - c\alpha)\varphi(0)$ on the natural
  merit suppresses it.
\item \textbf{The accuracy floor is the inner tolerance.} With damping, the
  outer natural residual descends monotonically to a floor that scales
  \emph{linearly} with the inner solver's tolerance; the finite-difference
  Jacobian does not bind (2nd- and 4th-order stencils produce identical
  iterates on the acceptance problems). Tightening the outer tolerance
  therefore means tightening \code{innerMagTol}, not refining the Jacobian.
\end{itemize}

\paragraph{Inexact inner solves (the forcing sequence).} There is no reason
to solve the linearized subproblem accurately while the outer iterate is
far from the solution: the classical inexact-Newton analysis of Dembo,
Eisenstat, and Steihaug \cite{dembo1982} shows convergence is retained when
each linearized system is solved only to a relative tolerance (a
\emph{forcing term}) that shrinks appropriately, and Eisenstat--Walker
\cite{eisenstat1996} give practical schedules tying the forcing term to the
current residual. \code{solveVI}'s factory overload implements the
residual-proportional member of that family,
$\mathit{innerTol}_k = \mathrm{clamp}(c \cdot \|r(z_k)\|^2,
\mathit{floor}, \mathit{cap})$ with $c = 10^{-2}$ --- exactly the
\code{min(5e-4, n0/100)} rule the project's reference Octave scripts have
always used, whose empirical signature (thousands of inner iterations on
the first linearizations, around a hundred on the last) was the original
motivation. Two library-specific notes: the tolerances are on \emph{squared}
residuals like everything else here, and the floor keeps the endgame
consistent with the fixed-tolerance floor law above (the outer residual
floor scales with the inner tolerance, so the forcing floor is what a
tight \code{outerTol} ultimately leans on).

The Jacobian itself is a 4th-order central difference with per-component
steps $h_j = \mathrm{stepRel}\cdot\max(|z_j|, 1)$ snapped to exactly
representable widths, $\mathrm{stepRel} = \varepsilon^{1/5}$ by default
(truncation error $O(\varepsilon^{4/5})$) --- the real-arithmetic
replacement for the reference implementation's complex-step
differentiation. The driver also carries a no-progress cutoff (stop, and
report \code{converged = false}, after $k$ consecutive outer iterations
without relative improvement, \emph{before} paying for the next Jacobian
and inner solve), added when a nonmonotone problem was observed freezing
the outer loop at a constant residual for twenty-plus iterations at a full
inner solve each (Chapter~\ref{ch:testing}).

\section{Mehrotra predictor-corrector for the mixed monotone LCP}
\label{sec:a-ipm}

\code{mehrotraIpm} is an infeasible primal-dual predictor-corrector
path-following method assembled from the standard recipes: the survey of
Potra and Wright \cite{potra2000} for the scheme, Gondzio \cite{gondzio2012}
for practical safeguards, and Gertz--Wright's OOQP \cite{gertz2003} for the
software pattern (including the layered linear-algebra interface of
Section~\ref{sec:a-flownewton}). For the mixed LCP over \K\ ---
$(Mz+q)_x = 0$, $0 \le y \perpc s = (Mz+q)_y \ge 0$ --- each iteration, at
strictly interior $(y, s)$:

\begin{enumerate}
\item forms the infeasibilities $r_x = (Mz+q)_x$, $r_s = (Mz+q)_y - s$ and
  the complementarity measure $\mu = y^\top s / m$;
\item factors the Newton matrix
  $\,K_\mu = M + \mathrm{blockdiag}\,(0,\ \mathrm{diag}(s/y))\,$
  \emph{once}, shared by both solves below;
\item \textbf{predictor}: solves for the affine-scaling direction
  (right-hand side built from $-r_x$, $-s - r_s$), computes the largest
  feasible step $\alpha_{\mathrm{aff}}$ and the resulting
  $\mu_{\mathrm{aff}}$, and sets the centering parameter
  $\sigma = \mathrm{clamp}\bigl((\mu_{\mathrm{aff}}/\mu)^3\bigr)$ ---
  Mehrotra's heuristic;
\item \textbf{corrector}: re-solves against the same factorization with the
  centered, second-order-corrected complementarity right-hand side
  $\sigma\mu\mathbf{1} - y \circ s - \Delta y_{\mathrm{aff}} \circ
  \Delta s_{\mathrm{aff}}$;
\item takes one \emph{common} damped step
  $\alpha = \min(1, \tau\,\alpha_{\max})$, $\tau = 0.995$, for all blocks.
  Separate primal/dual step lengths are an LP-only device; for LCP they
  break the convergence theory, so the implementation forbids them.
\end{enumerate}

The start is the data-scaled interior point $x = 0$,
$y = s = \beta\mathbf{1}$ with $\beta = \max(1, \sqrt{\|q\|_\infty})$
(OOQP's rule); no caller iterate can be used, hence no warm starts.
Termination and \code{residual} use the shared squared natural residual.
Three safeguards, each with its provenance: a \emph{sticky} free-block
regularization $\varepsilon_{\mathrm{reg}} I_n$ applied only after a Newton
solve returns non-finite values --- non-finite output, not singularity per
se, is the correct trigger, because a singular-but-consistent system can
solve finitely and needs no rescue (the discovery is recounted in
Chapter~\ref{ch:testing}); the fraction-to-boundary rule above;
and a stall return --- a damped step below $10^{-10}$ ends the run with
\code{converged = false} at the current iterate, since on data outside
the monotone theory (e.g.\ a nonmonotone linearization handed in by the
Josephy--Newton driver) blocked steps are an expected outcome, not an
error.

\paragraph{Why degeneracy does not hurt.} The central path
$\{z(\mu)\}_{\mu > 0}$ of a monotone LCP converges, as $\mu \to 0$, to the
\emph{analytic center of the optimal face}; path-following complexity
depends on problem data and $\log(1/\epsilon)$, not on any measure of how
flat or tied that face is, and superlinear behavior persists on degenerate
sufficient LCPs \cite{potra2009}. The growing conditioning of $K_\mu$ as
$\mu \to 0$ is benign for the computed steps in the sense analyzed by
Wright \cite{wright1997} and summarized by Gondzio \cite{gondzio2012}: the
correct response is the tolerance/iteration-cap termination, not a
conditioning guard. Empirically (Chapter~\ref{ch:testing}): 9/19/10
iterations on clean/degenerate/rank-deficient constructed LCPs, and 35--36
iterations at dimensions 10{,}838 and 14{,}601 on the banded flow problem
--- flat in both dimension and degeneracy, exactly as the theory predicts.

\section{The structured Newton solve for the flow LCP}
\label{sec:a-flownewton}

At scale, the IPM's cost is entirely its dense $O(d^3)$ factorization. The
flow LCP's structure admits an exact factory
(\code{makeFlowNewtonFactory}) that replaces it. Every formula in this
section was verified with the Maxima computer algebra system
\cite{maxima,macsyma}, symbolically and in exact rational arithmetic, by
the script \code{network/doc/ns2-newton-check.mac} --- ten checks: the
sink-block inverse, the KKT Jacobian against the assembled $M$,
recipe-equals-direct on three shapes including the $k_n = 1$ edge case, the
two-part Schur assembly, and symmetry --- all reproducible by running the
script. The prose here states the result and the two observations on which
the factory rests.

With the sink-major packing $z = [\,t \mid \mu \mid \lambda\,]$
(Section~\ref{sec:a-flowqp}), the IPM Newton matrix is
\[
K \;=\;
\begin{bmatrix}
Q_{\mathrm{blk}} + \mathrm{diag}(D_t) & E & d\\[2pt]
-E^\top & \mathrm{diag}(D_\mu) & 0\\[2pt]
-d^\top & 0 & D_\lambda
\end{bmatrix},
\]
where $\mathrm{diag}(D_t \mid D_\mu \mid D_\lambda)$ is the interior-point
complementarity diagonal $s/y$, $E$ scatters pairs to their sources
($E_{p\sigma} = 1$ iff pair $p$ ships from source $\sigma$), $d$ is the
pair-cost column, and $Q_{\mathrm{blk}} =
\mathrm{blockdiag}_n\,(Q_n \mathbf{1}\mathbf{1}^\top)$ --- each sink's
diagonal block is \emph{rank one}, with $Q_n = 2P_n/D_n^2$ the sink's
objective curvature.

\textbf{Observation 1 (per-sink Sherman--Morrison).} $W = (Q_{\mathrm{blk}}
+ \mathrm{diag}(D_t))^{-1}$ is block-diagonal by sink, and on sink $n$'s
slice
\[
W|_n \;=\; \mathrm{diag}(1/D_t) \;-\;
\frac{Q_n}{1 + Q_n \sigma_n}\, u_n u_n^\top,
\qquad u_n = (1/D_{t,p})_{p \in I_n},\quad
\sigma_n = \textstyle\sum_{p \in I_n} 1/D_{t,p},
\]
so applying $W$ costs $O(k_n)$ per sink --- linear in the number of kept
pairs overall.

\textbf{Observation 2 (the skew signs cancel; the Schur complement is
SPD).} Eliminating $t$ from the dual rows, the minus signs of the
$-E^\top, -d^\top$ blocks meet the minus signs inside the
back-substitution and cancel, leaving
\[
S \;=\; \mathrm{diag}(D_\mu \mid D_\lambda) \;+\; B^\top W B,
\qquad B = [\,E \;\ d\,],
\]
a symmetric positive-definite matrix of size
$(\text{numSources}+1)$ --- so Cholesky factorization applies even though
$K$ itself is unsymmetric. $S$ is assembled in two parts:
$B^\top \mathrm{diag}(1/D_t) B$ by an $O(\text{pairs})$ accumulation, minus
the per-sink rank-one corrections
$c_n (B^\top u_n)(B^\top u_n)^\top$ restricted to their support,
$O(k_n^2)$ each. The full solve per IPM iteration is then: apply $W$,
assemble and Cholesky-solve the small $S$, back-substitute --- about
$O(\sum_n k_n^2)$ work versus dense $d^3$, which at the 200-node keep-all
scale is the difference between $\sim 50$ seconds per iteration and
milliseconds, i.e.\ between a half-hour screened solve and a 4.5-second
exact one (Chapter~\ref{ch:testing}).

Correctness in floating point is asserted by \emph{backward-error} parity
bars against the dense factory,
$\|Kd - \mathit{rhs}\| / (\|K\|_\infty \|d\| + \|\mathit{rhs}\|) \le
10^{-12}$, always; direct step parity is asserted only on well-scaled
systems, since forward error carries $\mathrm{cond}(K)$ and the final
interior-point iterations make $K$ ill-conditioned by design. The
engine-side safeguard for \emph{any} external factory is the
\code{newtonCheckTol} consistency check of Section~\ref{sec:a-ipm}.

\section{The structured Newton solve for the fleet LCP}
\label{sec:a-fleetnewton}

The fleet optimizer (multiple asset and vehicle types;
\code{network/doc/fleet-formulation.md}, KKT system machine-checked by
\code{network/doc/fleet-mcp-check.mac}) generalizes the flow LCP, and its
Newton system admits the same treatment. Every formula in this section
was verified with Maxima \cite{maxima,macsyma}, symbolically and in exact
rational arithmetic, by the script
\code{network/doc/fleet-newton-check.mac} --- nine checks: the cell-block
Sherman--Morrison inverse; recipe-equals-direct-solve fully symbolically
on a one-cell, two-type shape (the two-column budget border the flow case
never exercises), in exact rationals on the three-cell fixture shared with
\code{fleet-mcp-check.mac} (so the two scripts pin the same variable
packing), and symbolically at the cell-size-one edge case; the $O(k)$
$W$-apply shortcut; the two-part Schur assembly against the dense
definition; and the symmetry of $S$, symbolically and in rationals ---
all reproducible by running the script.

The fleet variables are $z = [\,y \mid \mu \mid \lambda\,]$, packed
cell-major: one $y$ variable per (kept source--sink pair, capable vehicle
type), the variables of one \emph{demand cell} (sink, asset) contiguous;
one supply multiplier $\mu$ per (source, asset) cell; one budget
multiplier $\lambda_k$ per vehicle type. The interior-point Newton matrix
is
\[
K \;=\;
\begin{bmatrix}
Q_{\mathrm{blk}} + \mathrm{diag}(D_y) & E & R\\[2pt]
-E^\top & \mathrm{diag}(D_\mu) & 0\\[2pt]
-R^\top & 0 & \mathrm{diag}(D_\lambda)
\end{bmatrix},
\]
where $\mathrm{diag}(D_y \mid D_\mu \mid D_\lambda)$ is the
complementarity diagonal, $Q_{\mathrm{blk}} =
\mathrm{blockdiag}_c\,(Q_c \mathbf{1}\mathbf{1}^\top)$ has one
\emph{rank-one} block per demand cell $c$ with curvature
$Q_c = 2P_c/D_c^2$, $E$ scatters variables to their supply cells (one
unit entry per row), and $R$ scatters variables to their vehicle types
with weight $\rho_p$ (the round-trip miles per unit; one entry per row).
The two structural differences from Section~\ref{sec:a-flownewton}: the
budget border is $\mathit{numTypes}$ columns rather than one dense
column, and a cell mixes several sources \emph{and} several types.

Both observations of Section~\ref{sec:a-flownewton} carry over with cells
in place of sinks. The inverse
$W = (Q_{\mathrm{blk}} + \mathrm{diag}(D_y))^{-1}$ is block-diagonal by
cell with
\[
W|_c \;=\; \mathrm{diag}(1/D_y) \;-\;
\frac{Q_c}{1 + Q_c \sigma_c}\, u_c u_c^\top,
\qquad u_c = (1/D_{y,p})_{p \in c},\quad
\sigma_c = \textstyle\sum_{p \in c} 1/D_{y,p},
\]
applied in $O(k_c)$ per cell; and eliminating $y$ cancels the skew signs,
leaving the symmetric positive-definite Schur complement
\[
S \;=\; \mathrm{diag}(D_\mu \mid D_\lambda) \;+\; B^\top W B,
\qquad B = [\,E \;\ R\,],
\]
of size $\mathit{numSupplyCells} + \mathit{numTypes}$ (131 at the
70-node, $4 \times 3$ default; $\sim$1{,}200 at the production target),
factored by Cholesky. $S$ is assembled in two parts: each row of $B$ has
exactly two nonzeros ($1$ at the variable's supply cell, $\rho_p$ at its
type), so $B^\top \mathrm{diag}(1/D_y) B$ accumulates one $2 \times 2$
contribution per variable in $O(\mathit{numVars})$; then one rank-one
downdate $c_c\, (B^\top u_c)(B^\top u_c)^\top$ per cell, restricted to
the cell's support (its distinct sources plus its distinct types).

The remaining scale barrier after the factory was the explicit $M$
itself: $O(d^2)$ storage and residual products, terabytes at the
production target. The engine touches $M$ in exactly three computations
--- the field $Mz + q$, the \code{newtonCheckTol} drift guard, and the
built-in dense factory --- so \code{mehrotraIpm} accepts $M$ as an apply
operator (\code{MatrixApply}); the operator form requires an external
Newton factory (the built-in dense one needs the explicit matrix), and
the dense overload is proven unchanged bit for bit by an exact-equality
test. The fleet pipeline supplies the operator from its incidence lists
in $O(\mathit{numVars})$ --- one pass accumulating cell sums, one pass
emitting rows --- and never assembles $M$. The measured effect of the
factory and the operator together is the four-orders-of-magnitude
reduction of Chapter~\ref{ch:testing} (Table~\ref{tab:fleet}): the same 14{,}579-%
dimensional keep-all instance at 0.1\,s, identical iteration count and
objective, the residual differing from the dense-$M$ run only in the last
bit --- the expected signature of a changed summation order. Floating-%
point correctness carries the same bars as the flow factory: backward
error $\le 10^{-12}$ at any conditioning, direct parity on well-scaled
systems only, plus a parity case built from an \code{lcp} whose dense $M$
was never assembled.

\section{Semismooth Newton for the mixed NCP}
\label{sec:a-ssn}

\code{semismoothNewtonSolve} follows the line of De Luca--Facchinei--Kanzow
\cite{deluca2000} and the SEMI algorithm of Munson et al.\
\cite{munson2001}, with the penalized Fischer--Burmeister function of
Chen--Chen--Kanzow \cite{chen2000} as its default. Reformulate
\eqref{eq:mcp} as the nonsmooth square system
\[
\Phi(z) =
\begin{bmatrix} H(x, y)\\ \varphi\bigl(y_i,\, G_i(x,y)\bigr)_{i=1..m}
\end{bmatrix} = 0,
\qquad
\varphi_\lambda(a,b) = \lambda\,\varphi_{FB}(a,b)
+ (1-\lambda)\,a_+ b_+,
\]
with $\varphi_{FB}(a,b) = a + b - \sqrt{a^2+b^2}$, $a_+ = \max(0,a)$, and
$\lambda = 0.8$ (never below $0.5$: the penalty term then dominates the
merit and a merit-based stop can accept sign-violating points). The
penalization repairs plain FB's two documented weaknesses --- merit level
sets unbounded for merely monotone $F$, and near-flatness deep in the
positive orthant \cite{chen2000} --- and the library's own runs confirm
the comparison: on the deployment game the plain-FB variant failed soonest
of all configurations tried.

\paragraph{Generalized Jacobian.} Row $n+i$ of the Newton matrix is
$d_{a,i}\, e_{n+i}^\top + d_{b,i}\, (J_G)_i$ with the C-subdifferential
diagonals, away from the kink,
\[
d_a = \lambda\Bigl(1 - \frac{a}{\|(a,b)\|}\Bigr) + (1-\lambda)\, b_+ [a>0],
\qquad
d_b = \lambda\Bigl(1 - \frac{b}{\|(a,b)\|}\Bigr) + (1-\lambda)\, a_+ [b>0];
\]
at a kink ($a_i = b_i = 0$) the limiting-direction rule of
\cite{chen2000} (their Algorithm~4) selects a valid B-subdifferential
element from the direction $(z_i, (J_G z)_i)$ with $z$ the kink indicator.
One implementation subtlety is worth a permanent record: that indicator
is defined in the \emph{full} $z$-space (length $n+m$, ones at positions
$n+i$), not in $y$-space --- with $n = 0$ the two coincide, so pure-LCP
tests cannot catch the confusion; a mixed test hitting an exact kink can,
and did. Evaluation is stabilized per Munson et al.\ (their \S4.1): the
norm is computed overflow-safe by factoring out $|a|+|b|$, and for
$a+b > 0$ the cancellation-prone $a + b - \|(a,b)\|$ is replaced by the
algebraically equal $2ab/(a + b + \|(a,b)\|)$.

\paragraph{Globalization.} A three-tier direction ladder --- the Newton
direction if it satisfies the descent test
$\nabla\Psi^\top d \le -\rho\|d\|^{p}$ ($p > 2$, the De Luca et al.\
condition); else a Levenberg--Marquardt direction with damping
$\lambda_{LM} \propto \|\Phi\|^2$ (the error-bound rule that retains fast
local convergence on non-isolated solution sets, per the H\"older
subregularity analysis of \cite{ahookhosh2019}, with no nonsingularity
assumption); else the negative gradient --- followed by a
\emph{directional} Armijo search on $\Psi = \tfrac12\|\Phi\|^2$ against a
trailing-window baseline (\code{nonmonotoneMemory}; window 1 is the citable
monotone theory, window 4 is SEMI's production value). The middle tier is
necessary: singular Newton matrices are \emph{common} on degenerate
problems \cite{munson2001}, and the gradient tier alone is unreliable.

\paragraph{Contract-level decisions.} Iterates are not confined to $K$;
a trial point where the model throws (a domain violation, e.g.\ crossing a
pole of a rational vector field) counts as merit $+\infty$ \emph{inside
the line search only} --- the step backtracks away --- while non-finite
values at accepted points remain hard errors. Termination is on the
library's squared natural residual rather than on $\Psi$: the
$\Psi$-stop under-enforces sign conditions for small $\lambda$ (a
documented pitfall), and the natural residual keeps the engine comparable
with every other. The returned iterate is the \emph{best visited} in that
measure: the nonmonotone search may deliberately climb above its best point
and then stall before recovering it --- observed concretely (a run
returning squared residual $1.7\times 10^{4}$ after visiting
$8.2\times 10^{3}$) before the semantics were fixed.

\section{From the GMS Model statement to the MCP}
\label{sec:a-gmsmcp}

The front end of Chapter~\ref{ch:manual} builds a \code{VIModel} from a
parsed \code{Model} statement\footnote{``GAMS'' is a registered trademark
of GAMS Development Corporation; the front end parses an incompatible
subset of the GMS language and is not endorsed or certified by GAMS
Development Corporation.}; this section states the construction and the
two arguments it rests on. A \code{Model} statement is a list of pairs
$e_1.v_1, \ldots, e_k.v_k$, each an equation family and a variable
family. The construction is:

\begin{itemize}
\item a pair whose variable is \emph{free} contributes its rows to $H$,
  and the relation must be \code{=e=};
\item a pair whose variable is \emph{positive} contributes its rows to
  $G$, whatever the relation (\code{=g=} or \code{=e=});
\item every row's residual is $\mathit{lhs} - \mathit{rhs}$, which
  matches the library's orientation directly: the GMS pairing
  ``$\mathit{body} \geq 0$ complementary with $v \geq 0$'' is exactly
  $0 \leq G \perp y \geq 0$.
\end{itemize}

The block assignment is decided by the \emph{variable's kind}, with the
relation reduced to a consistency check, because of the
mixed-complementarity reading of an equality paired with a positive
variable: $0 \leq v \perp g(z) \geq 0$ holds wherever the equality
$g(z) = 0$ holds with $v \geq 0$, and conversely, whenever the model
guarantees $v$ strictly positive at any solution, complementarity forces
$g(z) = 0$ and the equality is recovered exactly. The influence game is
the corpus case: its intermediate variables are declared
\code{Positive} but are defined by equalities of the form
$\varepsilon$ plus a sum of nonnegative terms, hence strictly positive,
so the complementarity reading and the equality reading coincide (the
same interiority argument as step 3 of the hand recipe in
Chapter~\ref{ch:manual}).

Each family expands over its domain in row-major order, and the
construction requires the equation definition's domain and the paired
variable's domain to resolve, dimension by dimension, to the same base
sets; element $t$ of the expanded equation family then pairs with element
$t$ of the expanded variable family, and the packed position of a
variable element equals the row position of its equation --- one identity
the tests re-verify on every model. The start $z_0$ packs the \code{.L}
levels in the same order. Upper bounds (\code{.UP}) are carried on the
result but not imposed: $K$ remains the orthant, and each corpus model
that sets bounds also carries the argument that its pool constraints
imply them at any solution.

When every equation in the model is affine, $F(z) = Mz + q$ with a
constant $M$, and the exact data are recoverable from evaluations alone:
$q = F(0)$ and $M e_j = F(e_j) - q$, so $\mathrm{dim}+1$ evaluations
assemble $(M, q)$. Affinity itself is checked at a random point $z$ by
comparing $F(z)$ against $Mz + q$: for the censused expression grammar a
nonaffine $F$ differs from its affine interpolant except on a
lower-dimensional set, so a random draw detects nonaffinity with
probability one (and the check is cheap either way). On the logistics
corpus model the check agreed to a relative $2.6\times 10^{-15}$ and the
assembled LCP went to the interior-point engine of
Section~\ref{sec:a-ipm}; Chapter~\ref{ch:testing} records the
measurements.

\section{The application problems}
\label{sec:a-apps}

The library's three almost-real-world problems are treated here as peers,
one subsection each: the quadratic flow-planning problem (a single-planner
convex optimization), and two \emph{games} --- the strategic allocation of
effort (with a risk-averse variant forthcoming) and the
pre-position-and-deploy family --- whose Nash equilibria are posed directly
as MCPs. The games arrive as GAMS models with verified
solutions and are transcribed one-to-one into \code{VIModel}s (the recipe
is Chapter~\ref{ch:manual}); each subsection records what an OR reader
needs to audit the formulation.

\subsection{The quadratic flow-planning problem}
\label{sec:a-flowqp}

The statements of this subsection are established in the project's working
documents (\code{network/doc/formulation.md}, \code{reduction.md}) with
full proofs; they are transcribed here in summary, not re-derived.

\paragraph{Formulation.} On a node set $V$ with supply caps $C_i$, demands
$D_i$, priorities $P_i > 0$, positive arc costs $c_{ij}$ (asymmetric;
positive diagonal), and ton-mile budget $L$: choose injections
$S_i \in [0, C_i]$, resupplies $R_i \ge 0$, and flows $f_{ij} \ge 0$ to
minimize the dimensionless weighted shortfall
\[
\theta(R) = \sum_{i \in V_D} P_i \Bigl(\frac{D_i - R_i}{D_i}\Bigr)^{2},
\qquad V_D = \{\,i : D_i > 0\,\},
\]
subject to flow balance, the budget
$\sum_{ij} c_{ij} f_{ij} \le L$, and the \emph{delivery} constraint
$R_j \le \sum_i f_{ij}$. The delivery row is a spec-audit finding
(``F1''): without it, the written balance lets a node self-supply for free,
contradicting the stated intent; adding it is equivalent to splitting each
node into supply/delivery ports with a nonnegative through-flow, and forces
self-supply through $f_{jj}$ at cost $c_{jj}$. The problem is a linearly
constrained convex QP --- strictly convex in $R_{V_D}$ (diagonal Hessian
$2P_i/D_i^2$) --- feasible and compact, so a minimizer exists and the
optimal $R^\ast$ is unique (the optimal routing need not be).

\paragraph{Reduction (Lemma R1).} With $\hat d_{mn}$ the arc-respecting
shortest-path cost from source $m$ to sink $n$ (the diagonal corrected so a
self-supplied ton uses the self-arc or a round trip), the pair problem
\[
(P_\varepsilon)\quad
\min_{t \ge 0}\ \theta\bigl(R(t)\bigr)
+ \varepsilon \sum_{m,n} \hat d_{mn} t_{mn},
\qquad
R_n(t) = \sum_m t_{mn},\quad
\sum_n t_{mn} \le C_m,\quad
\sum_{m,n} \hat d_{mn} t_{mn} \le L,
\]
has (at $\varepsilon = 0$) the same optimal value as the full problem, and
any feasible $t$ expands --- along stored shortest routes --- to a feasible
full plan with the same $R$ and no greater budget use. The proof is by flow
decomposition on the port-split digraph; the ``expand'' direction is the
production unpacker. The reduction collapses $M^2 + 2M$ variables to
$|V_C||V_D|$ (at 200 nodes: $\sim$40{,}400 to $\le$10{,}000).

\paragraph{KKT $\leftrightarrow$ LCP.} Substituting $R = R(t)$ out, the KKT
system of $(P_\varepsilon)$ in the variables
$z = [\,t \mid \mu \mid \lambda\,]$ (pair flows; capacity duals; budget
dual) is a \emph{pure} NCP over the nonnegative orthant --- no free block
--- whose affine map has positive semidefinite $M$: the objective
contributes the sink-blocked rank-one curvature $Q_n \mathbf{1}\mathbf{1}^\top$
($Q_n = 2P_n/D_n^2$), and the constraint rows enter skew-symmetrically,
contributing nothing to the symmetric part. Monotonicity is therefore
structural, and every monotone-theory engine applies. This packing (each
sink's pairs contiguous) is also exactly what the structured factory of
Section~\ref{sec:a-flownewton} exploits.

\paragraph{Screening and the certificate (Proposition R3).} Restricting to
a screened pair set and solving, optimality of the padded solution for the
\emph{unrestricted} problem is certified by dual feasibility of every
excluded column:
\[
\frac{2P_n}{D_n}\cdot\frac{D_n - R_n^\star}{D_n}
\;\le\; (\lambda^\star + \varepsilon)\,\hat d_{mn} + \mu_m^\star
\qquad \text{for every excluded } (m,n);
\]
violated pairs are added and the problem re-solved --- column generation
with a convex master, finitely terminating at the exact optimum. A
\code{certificateSlack} loosens the check to a \emph{bounded-suboptimality}
certificate. The practical failure mode at scale is not incorrectness but
abrupt growth of the kept set: on near-tied geometries each certificate
round adds many near-zero-reduced-cost pairs, each round at full solve
cost (one observed round grew the kept set from $\sim$1{,}900 to 10{,}737
of 14{,}500). The structured keep-all solve removes the screen, and with
it the certificate loop, entirely: nothing excluded, nothing to certify.

\paragraph{Tie-break (Proposition R4).} The linear term
$\varepsilon \sum \hat d_{mn} t_{mn}$ selects minimal-ton-mile plans among
near-optimal routings (the observed flow degeneracy is exactly ties in this
direction) at a certified objective cost
$\theta(t^\varepsilon) \le \theta^\star + \varepsilon L$; the shipped
$\varepsilon = 10^{-8} \sum_n P_n / L$ prices the tie-break at $10^{-8}$ of
full-scale shortfall.

\paragraph{Nondimensionalization.} The solver works in scaled units ---
tons divided by the largest demand, miles by the largest cost --- an
\emph{exact} change of variables under which the optimum is invariant. The
rescaling has practical effect: the projection engines' effective
conditioning depends on the raw scales, and the lane-oracle benchmark
improved from 25\,s to 0.04\,s on nondimensionalization alone. All
tolerances are expressed in scaled units; results are returned in real
units.

\paragraph{The fleet generalization: a conservative formulation.} The
fleet model (\code{network/doc/fleet-formulation.md}) extends the flow
QP to $|A|$ asset types and $|K|$ vehicle types. Its full form couples
the assets on every link through two aggregate constraints --- cargo
weight $\sum_a w_a x^a_{ij} \le \sum_k T_k u^k_{ij}$ and cargo area
$\sum_a s_a x^a_{ij} \le \sum_k A_k u^k_{ij}$, which jointly permit
\emph{mixed loading} (one vehicle carrying a weight-bound and an
area-bound asset simultaneously) --- and routes the vehicles $u^k$ as a
per-type circulation. That problem has $(|A| + |K|)M^2 + 2|A|M$
variables, roughly 35{,}000 at the 70-node default, and its KKT matrix
is far beyond the dense factorizations every current engine relies on
(\code{bsHe94b} factors $(M+I)$ once; the interior-point engine factors
a Newton matrix each iteration). The shipped fleet optimizer therefore
solves a \emph{conservative} formulation instead: each asset is loaded
at the per-type rate $\kappa_{ak} = \min(T_k/w_a,\, A_k/s_a)$ units per
vehicle, which satisfies both aggregate constraints without exploiting
mixed loading, and vehicles travel out-and-back along shortest routes,
which satisfies the circulation constructively. Under these two
conventions --- the same ones the greedy fleet planner uses --- the
vehicle variables reduce to $|K|$ budget rows over per-(asset, type)
mileage rates, the per-asset shortest-route reduction applies, and the
resulting QP has the same monotone pure-NCP structure as
$(P_\varepsilon)$ with rank-one curvature blocks per demand cell and
$|K|$ skew budget borders in place of one. The conservative optimum is
an upper bound on the full model's optimum, and it is the exact optimum
of the model the greedy planner bounds from above, so the validation
sandwich $\theta_{\text{ration}} \le \theta^\ast \le
\theta_{\text{greedy}}$ compares like with like. Solving the full
link-coupled QP requires sparse solver machinery --- sparse assembly and
factorization, or a matrix-free engine --- that the library does not yet
have; it is recorded as future work (\code{network/plan.md}, G5i).

\subsection{Strategic allocation of effort (risk-neutral)}
\label{sec:a-saoe}

The general problem: $M$ actors with weights $w_a$ allocate effort
$e_{ap} \ge 0$ across $N$ options under budgets $\sum_p e_{ap} \le w_a$;
option strengths $\sigma_p$ aggregate efforts, option probabilities are
$\sigma_p / \gamma$ with $\gamma = \sum_p \sigma_p$, and each actor's
stationarity row prices its marginal influence on its expected reward
against a shadow price $\beta_a$ on its budget. The Nash equilibrium is
posed directly as an MCP. A mixing parameter $\alpha$ interpolates the
strength map between purely linear ($\alpha = 1$) and purely quadratic
($\alpha = 0$); at $\alpha = 1$, the setting of the test instance, the only
nonlinearity left is rational: the $\gamma^2$ denominator in the
stationarity rows. The concrete instance in the acceptance suite
(\code{alloceff01cm.gms}, transcribed in
\code{test/gams\_alloceff\_test.cpp}) is one example of this family, at
$M = 6$ actors and $N = 10$ options.

\emph{Coalition persistence.} A feature of this problem --- in its applied
setting, the \emph{policy assemblage problem} --- observed consistently
in the parametric runs: as the risk-aversion parameter sweeps, the
\emph{coalition structure} (which actors pool their effort on a common
option) persists across wide parameter ranges and even across equilibrium
switches, while the option \emph{labels} are labile --- the same multiset
of option probabilities recurs with the probabilities attached to different
options. Interests, like economic interests generally, carry natural
positive and negative correlations, and it is those correlations --- not
the option identities --- that the equilibria appear to organize around.
The model's author has seen the same behavior in other coalition-formation
problems; we record it here as a feature of \emph{this} problem and
deliberately make no stronger claim: establishing it as a regularity would
take sweeps over thousands of instances (uniform and variously structured
payoffs), an empirical characterization, and an algebraic delineation of
when it does and does not occur.

\emph{Provenance and interior equilibria.} The test instance is not
arbitrary: it is the ``small example'' of the PME framework paper (its
\S4.4.6), whose rewards are value-added changes of the six sectors of a
synthetic Leontief input-output economy under ten approximately
GDP-neutral tax policies --- which is why its interest correlations are
economically structured rather than random. The paper's own solve of this
instance (forward-backward splitting, Section~\ref{sec:a-fbs}) reported an
\emph{interior} point --- supports $\{4, 6, 9\}$ with actor~1
\emph{splitting} its effort across options 4 and 6, the paper's
``hedging'' narrative --- distinct from the all-in vertex equilibrium $E$.
Whether that point is a genuine second (interior) equilibrium or an
artifact of the historical run's looser, step-based stopping rule
($\sim$15{,}000 iterations to $\sim$$10^{-6}$) is an \emph{open question}:
this library's faithful reimplementation of the same splitting method
converges to the vertex $E$ from the same start --- at both the historical
strength floor $\epsilon = 0.1$ and the library default $0.0108$, which
also makes $E$'s identity robust to a tenfold change in $\epsilon$ and
gives $E$ a fourth independent method-class confirmation. The question is
decidable by restarting the iteration from the historical allocation and
measuring its natural residual; it is recorded here rather than resolved.

The transcription makes two formally-justified moves. First, the GAMS
intermediates ($\mathrm{nfv}_p$, $\sigma_p$, $\gamma$) are declared
\code{Positive} in the source but are treated as \emph{free} variables with
their defining equalities as $H$ rows: at any point satisfying those
equalities, $\mathrm{nfv}_p \ge \epsilon > 0$ (an explicit floor plus a sum
of nonnegatives), hence $\sigma_p > 0$ and $\gamma > 0$, so the bounds
never bind and the KKT systems coincide. Second, the individual box bounds
$e_{ap} \le w_a$ are dropped as implied by the budget row; at the corner
where one effort exhausts the budget, the budget's multiplier absorbs the
box multiplier in the same stationarity row, so the primal solution sets
coincide. Empirical closure: two independent engine classes (semismooth
Newton; Josephy--Newton over the interior-point inner) converge to the same
equilibrium, which the author confirmed is the one PATH reaches --- also
the reference equilibrium of the library's independent, reduced-form SAOE
implementation. A sign error essentially anywhere in the transcription
would have split those three.

\subsection{Strategic allocation of effort with risk aversion}
\label{sec:a-saoe-risk}

\emph{[Reserved: a variant of the preceding problem in which the actors are
somewhat risk-averse. The problem statement is forthcoming and will be
inserted here.]}

\subsection{Pre-position and deploy (PPD)}
\label{sec:a-ppd}

PPD is a \emph{family} of two-sided pre-position/deploy games, developed in
versions of increasing sophistication (several more versions beyond the one
tested here are in preparation). The common core: each side pre-positions
forces at bases, assigns a mixed strategy over postures, and ships flows
from bases to contested locations under logistical (ton-mile) budgets;
payoffs are shares of location values determined by the forces that arrive;
and each side's stationarity, paired with its decision variables, poses the
Nash equilibrium directly as an MCP. Later versions enrich the interaction
--- interdiction along routes, escort and attacker allocations, staged
combat --- while preserving that MCP shape, so each version drops into the
library as data plus a \code{VIModel} builder and joins the same
multi-engine acceptance harness.

The family's defining mathematical character comes from its smooth
attrition modeling. Combat uses survivor functions of the form
$h(x) = x^2/(x + c)$, $c > 0$, with three consequences that shape solver
behavior across every version: $h(x) \le x$ (survivors never exceed forces
committed, an algebraic identity); $h'(0) = 0$ with $h'' > 0$ near zero
(small commitments are destroyed, mass is rewarded), so interior starts
are \emph{essential} --- a route at zero shows no first-order reason to
open, zero-flow local equilibria exist (economy of force in smooth form),
and the GAMS \code{.L} levels are part of the model, transcribed exactly;
and the composite payoff has mixed curvature --- concave in own variables
from strong positions, convex from weak ones (``mass-or-abstain'') --- so
the MCP is nonmonotone \emph{by design} and multiple Nash equilibria are
expected. The $\epsilon$-guards on every denominator place poles of the
vector field just outside the nonnegative orthant, a fact that matters for
solver behavior (Section~\ref{sec:a-chain}).

Versions in the acceptance suite to date:

\begin{itemize}
\item \textbf{\code{deploy\_v07.gms}} (transcribed in
  \code{test/gams\_deploy\_test.cpp}): movers, attackers, and defenders
  with interdiction along the routes and ratio combat throughout. Each side
  pre-positions movers, mixes over five postures, ships to four contested
  locations, and allocates escorts to its own flows and attackers against
  the opponent's routes (escorts and attackers are pooled per strategy,
  with no ton-mile cost). Combat is two smooth ratio stages per route ---
  attackers surviving the escorts,
  $\mathrm{BA2} = \mathrm{BA}^2/(\mathrm{BA} + \mathrm{RE} + \epsilon)$,
  then movers surviving the remaining attackers,
  $\hat f = f^2/(f + \mathrm{BA2} + \epsilon)$ --- followed by force-ratio
  shares of location values at each node. The stationarity rows carry the
  \emph{exact} chain-rule gradients of the two-stage survivor composition,
  verified symbolically by the model's author with Maxima
  \cite{maxima,macsyma} (\code{v07\_check.mac}) and transcribed literally,
  every $\epsilon$-guard included. Dimensions: 4 free rows (two
  probability-simplex and two pre-positioning equality multipliers) and 446
  complementarity rows. Two Nash equilibria are currently verified,
  discovered by the library itself and recorded in the acceptance test as
  a roster: both share the strategy supports, and they form a mirror-image
  pair in which exactly one side's per-strategy probability cap $\rho$ is
  active (Red's in one, Blue's in the other) --- suggestive of a symmetry
  family, recorded as an observation rather than a claim of enumeration.
\end{itemize}

\section{The alternating chain, formally}
\label{sec:a-chain}

No engine in this library carries a global convergence theory for
nonmonotone MCPs, and on the deployment game every standalone engine fails
--- each in its own way. The alternating chain
(\code{alternatingChainSolve}) is an engineered composition whose claims
are deliberately modest and checkable. We state them.

\paragraph{Construction.} Fix a globalizer $\mathcal{G}$ and finisher
$\mathcal{F}$ (stage solvers returning \code{VIResult}s), the projector
$P_K$, and the residual functional $\rho(z) = \|z - P_K(z - F(z))\|^2$
evaluated by the chain itself. A \emph{round} from incumbent best $b$ is
\[
z_0 = P_K(b), \qquad
z_1 = \mathcal{G}(z_0), \qquad
z_2 = \mathcal{F}(z_1),
\]
with the best point updated against $\rho$ after each stage. A stage that
\emph{throws} is treated as a stalled stage: the round continues from the
last good point (on a globalizer failure the finisher starts from $z_0$),
and nothing is lost --- the incumbent is kept. Rounds repeat up to
\code{roundsMax}; a round that fails to improve the incumbent
(\code{improveFactor}, default: any strict improvement continues) triggers
either termination or, with \code{perturbScale} $= s > 0$, a
\emph{perturbed restart}: the next round starts from
$P_K\bigl(b + \delta\bigr)$, $\delta_i \sim U[-1,1]\cdot s\cdot
\sqrt{\rho(b)}$ from a fixed-seed generator.

\paragraph{What is guaranteed.} (i) \emph{Boundedness and exact
reporting}: at most \code{roundsMax} rounds of internally-bounded stages;
the returned point is the best visited under the fixed functional $\rho$,
with \code{converged} true iff $\rho(\mathrm{best}) < \code{magTol}$; the
whole procedure is a deterministic function of its inputs (the
perturbation seed is fixed), so runs are reproducible. (ii) \emph{The
stagnation lemma} (elementary but essential): both stages being
deterministic, a round that starts from $P_K(b)$ and fails to improve $b$
will, if restarted verbatim, repeat identically --- so no stopping-rule
refinement can rescue a stagnant chain, and \emph{perturbation is the only
within-chain remedy}. The empirical form was exact: two consecutive runs
of a stagnant configuration were digit-for-digit identical. (iii)
\emph{Perturbation does not obstruct convergence}: $\delta$ scales with
$\sqrt{\rho(b)}$, so it vanishes as the chain approaches a solution.

\paragraph{What is argued, not guaranteed.} The composition targets
\emph{complementary failure sets}. The Josephy--Newton/IPM globalizer fails
where the current linearization is indefinite --- a local property of
$J(z_k)$; the semismooth finisher fails at merit-geometry obstructions ---
in this problem class, characteristically \emph{outside} $K$ beside the
$\epsilon$-guard poles, where its line search fails against infinite merit
values (the residual breakdown at one observed stall showed negative
escort allocations of $-0.34$ against $\epsilon = 0.01$), or at
nonmonotone merit plateaus. Neither failure property implies the other at
a handoff point, and the inter-round projection removes the finisher's one
systematic way of degrading the handoff (leaving $K$ toward the poles). A
deploy-specific check closes a potential hole in that argument: escorts
projected \emph{to} zero retain a first-order re-entry signal
($\partial \hat f/\partial \mathrm{RE} > 0$ at $\mathrm{RE} = 0$ on a
contested flowing route), unlike flows at zero, so the projection cannot
create new $h'(0) = 0$ dead zones. There is no global theorem here and we
claim none; the evidence is Chapter~\ref{ch:testing}'s case study, whose
endpoint is a machine-precision solution ($\rho \approx 10^{-9}$ to
$10^{-15}$) of a 450-dimensional nonmonotone MCP that no component engine
solves alone.

\paragraph{Precedents.} The design consciously follows PATH's practice ---
crash phase, restarts, perturbation on failure \cite{ferris1998} --- and,
one level down, the library's own affine chain (Solodov--Svaiter
globalizing He's contraction). The perturbed restarts also serve as an
equilibrium \emph{sampler}: distinct trajectories have already produced two
verified equilibria, and the acceptance test is built as a growing roster
for exactly that reason.

\section{Measurement conventions}
\label{sec:a-measure}

For auditing the numbers in this report:

\begin{itemize}
\item \textbf{Residuals are squared Euclidean norms} of the natural map
  \eqref{eq:natres}, in the problem's working units (for the flow pipeline:
  nondimensionalized units). ``Converged at $10^{-8}$'' means
  $\|r\|^2 < 10^{-8}$, i.e.\ $\|r\| < 10^{-4}$.
\item \textbf{Iteration counts are not comparable across engine classes.}
  One \code{bsHe94b} iteration is a back-substitution; one \code{mehrotraIpm}
  iteration is a full factorization (its \code{iterMax} is set in hundreds
  where a projection engine's is in hundreds of thousands); one
  \code{solveVI} outer iteration contains a complete inner solve plus a
  finite-difference Jacobian ($4(n+m)$ evaluations of $F$); composite
  results carry the split explicitly (\code{iter} outer,
  \code{innerIters} total inner).
\item \textbf{Wall-time attribution}: for factorization-bound engines,
  count factorizations times $O(d^3)$ (or the structured cost of
  Section~\ref{sec:a-flownewton}); for the Josephy--Newton driver on
  black-box models, the FD Jacobian's $4(n+m)$ model evaluations per outer
  step frequently dominate; for projection engines, iteration count is the
  cost, with factor-once vs.\ factor-every-iteration separating He from
  Han.
\end{itemize}

The annotated bibliography at the end of this report lists every source
cited here; local copies are kept under \code{doc/}, and each entry's
annotation states what the library took from it.

% ----------------------------------------------
% Copyright Ben Paul Wise. All Rights Reserved.
% ----------------------------------------------
